% !TeX root = ../main.tex
\section{Conclusion}\label{sec:conclusion}

At QD=1 the Micron 2210's write buffer makes \qty{4}{\kibi\byte} random writes faster than reads. Block size moves the
bottleneck from per-request latency to bandwidth, but even at QD=32 sequential reads reach only
\qty{56}{\percent} of the PCIe link, so the NAND back end is the ceiling, and writes saturate the
\ensuremath{\sim}\qty{1.05}{\gibi\byte\per\second} SLC ingest rate even at QD=1. Mixing reads and writes
is worse than either pure workload. The read knee lies at QD~\csvq{queue_depth_knee}{mode=read}{knee_qd}
under an explicit \qty{10}{\percent}-per-doubling rule; beyond it latency doubles per doubling of QD, and
$L/\text{QD}$ of 0.987--0.999 confirms that the saturation is the device's, not the harness's.

Device state matters as much as configuration. On a freshly trimmed, idle drive the SLC cache held
\csvval[1]{sustained_summary}{cliff_gib}\,GiB, against \qty{3.8}{} and \qty{3.1}{\gibi\byte} in earlier
runs, plausibly because it scales with free blocks; past it writes fell $\ssdburstratio\times$ and did
not recover after a minute of idle time. The same exhaustion briefly cut writes to the QLC rate in the
block-size and queue-depth sweeps, and it made the write knee rule return a meaningless QD=2. Reads did
not slow with span beyond \qty{4}{\gibi\byte}, but the write history of the data made
\enquote{sequential} reads 6--9$\times$ slower, and in our first run recent heavy writes stretched p99.9/p50 from
\ssdtailratio{128}{p999}$\times$ to $24\times$. Temperature stayed far below the warning threshold.

The sharpest lessons are methodological. Without \code{O_DIRECT} the same workload reported page-cache
fiction, and the first dataset was taken on the HDD until the \code{rotational} flag gave it away. A
storage measurement can be trusted only if it bypasses the cache, verifies its effective queue depth,
controls the data's write history and the drive's recent workload, and confirms which device it is
talking to.
